%!TEX program = lualatex
% SPDX-License-Identifier: MIT
% Copyright (c) 2026 the repository contributors
% Copy this ONE file to start a talk. It needs only the Atlas theme in sty/.
% In this repo: ./build-beamer.sh --starter [--plain] [--handout]
% Elsewhere: TEXINPUTS=/path/to/latex-atlas/sty: lualatex talk.tex (three times).
\documentclass[aspectratio=169,11pt]{beamer}
\usetheme{AtlasGrand} % Or \usetheme[plain]{AtlasGrand} for black on white.

% EDIT: your metadata. Keep the title short enough for the opening layout.
\title{One idea, clearly seen}
\subtitle{A claim, its evidence, and what follows}
\author{Your name}
\institute{Your institution}
\date{\today}

% Local layout choices. Group their use in {...} to restore the theme.
\newcommand{\TalkQuiet}{%
  \setbeamertemplate{background canvas}{\color{AtlasPaper}\vrule width\paperwidth height\paperheight}%
  \setbeamertemplate{footline}{}%
}
\newcommand{\TalkWhite}{%
  \setbeamertemplate{background canvas}{\color{white}\vrule width\paperwidth height\paperheight}%
  \setbeamertemplate{footline}{}%
}

% The example chart is illustrative, not a research result. Both views
% use the SAME drawing; the detail only changes the clipping rectangle.
% Delete this macro when replacing the two figure commands below.
\newcommand{\ExampleChart}[1]{%
  \begin{tikzpicture}[x=10mm,y=7mm,font=\small]
    \ifnum#1=0\def\leftedge{-.9}\else\def\leftedge{4.5}\fi
    \begin{scope}
      \clip (\leftedge,-.8) rectangle (8.5,5.3);
      \fill[white] (\leftedge,-.8) rectangle (8.5,5.3);
      \foreach \v in {0,2,4,6,8,10}{
        \draw[AtlasGold!25!white,line width=.3pt] (0,\v*.45)--(8.2,\v*.45);
        \node[anchor=east,text=AtlasInk] at (-.1,\v*.45) {\v};
      }
      \foreach \x/\label/\v in {1.5/A/4,3.5/B/6,5.5/C/5,7.5/D/9}{
        \filldraw[fill=AtlasGold!35!white,draw=AtlasInk,line width=.5pt]
          (\x-.45,0) rectangle (\x+.45,\v*.45);
        \node[above,text=AtlasInk] at (\x,\v*.45) {\v};
        \node[below,text=AtlasInk] at (\x,-.12) {\label};
      }
      \draw[AtlasInk] (0,0)--(8.2,0);
    \end{scope}
    \pgfresetboundingbox
    \path[use as bounding box] (\leftedge,-.8) rectangle (8.5,5.3);
  \end{tikzpicture}%
}
% EDIT: substitute your own figure and detail. Size each within its slot.
% Example: \includegraphics[width=148mm,height=65mm,keepaspectratio]{figure.pdf}
\newcommand{\TalkFigure}{\resizebox{!}{62mm}{\ExampleChart{0}}}
% Prefer a separately exported panel that retains the axes and legend it needs.
% Example: \includegraphics[width=57mm,height=50mm,keepaspectratio]{detail.pdf}
\newcommand{\TalkDetail}{\resizebox{!}{50mm}{\ExampleChart{1}}}

\begin{document}

% 1. OPENING — metadata, generous space, and the dark palette.
{\AtlasGrandNight
\begin{frame}[plain,noframenumbering]
  \centering
  {\color{AtlasGoldLight}\pgfornament[width=26mm,color=AtlasGoldLight]{71}\par}
  \bigskip
  {\color{AtlasIvory}\fontsize{32}{36}\selectfont\scshape\inserttitle\par}
  \medskip
  {\color{AtlasIvory}\large\itshape\insertsubtitle\par}
  \bigskip\bigskip
  {\color{AtlasGoldLight}\insertauthor\par\smallskip
    \small\insertinstitute\par\insertdate\par}
\end{frame}}

% 2. ORDINARY CONTENT — native blocks and speaker-paced reveals.
% Handout mode displays both items together. Remove [<+->] for a static slide.
\section{The claim}
\begin{frame}{Make one comparison worth remembering}
  \begin{block}{The claim}
    The final group stands apart from the others.
  \end{block}
  \medskip
  \begin{itemize}[<+->]
    \item Show the full comparison before drawing attention to a detail.
    \medskip
    \item Explain what the difference means, and what it leaves unresolved.
  \end{itemize}
  \note{Replace this example claim with yours. Keep the conditions and limits
    of the claim close to the evidence. This sample uses illustrative values.}
\end{frame}

% 3. FULL FIGURE — a neutral field preserves source colors and annotations.
% Plain frames still count as frames unless noframenumbering is specified.
{\TalkWhite
\begin{frame}[plain]
  \begin{tikzpicture}[remember picture,overlay]
    \node[anchor=north west,text=AtlasInk,font=\large\scshape] at
      ([xshift=8mm,yshift=-5mm]current page.north west) {The full comparison};
    \node at ([yshift=-1mm]current page.center) {\TalkFigure};
    \node[anchor=south,text=AtlasInk,text width=143mm,align=center,font=\small] at
      ([yshift=5mm]current page.south) {Illustrative data\quad /\quad Four groups, one common scale.};
  \end{tikzpicture}
\end{frame}}

% 4. ENLARGED DETAIL — use a useful crop, then state its interpretation.
\section{The evidence}
\begin{frame}{The last pair makes the difference visible}
  \begin{columns}[c,onlytextwidth]
    \begin{column}{.46\textwidth}\centering\TalkDetail\end{column}
    \begin{column}{.48\textwidth}
      \AtlasGrandKicker{Read the detail}
      {\color{AtlasRubric}\fontsize{30}{34}\selectfont $5\longrightarrow9$\par}
      \medskip
      Group D is 80\% above group C in this illustrative comparison.\par\medskip
      \AtlasGrandScholium{The next question}{What mechanism could explain
        the difference? What evidence would distinguish it from alternatives?}
    \end{column}
  \end{columns}
\end{frame}

% 5. QUIET QUOTATION — EB Garamond Italic, 27pt with 31pt leading.
% Add the actual attribution below when quoting someone else's words.
{\TalkQuiet
\begin{frame}[plain]
  \centering
  {\fontsize{27}{31}\selectfont\itshape
    Give the evidence room.\\
    Then give the conclusion a voice.\par}
  \bigskip\bigskip
  {\small\color{AtlasGold}\textsc{A pause before the conclusion}\par}
\end{frame}}

% 6. CLOSING — your takeaway and a useful next step or discussion question.
{\AtlasGrandNight
\begin{frame}[plain]
  \centering
  {\color{AtlasGoldLight}\small\textls[140]{\textsc{The takeaway}}\par}
  \bigskip
  {\color{AtlasIvory}\fontsize{30}{35}\selectfont\scshape
    Make the comparison count.\par}
  \bigskip
  {\color{AtlasIvory}\large One finding. One implication. One next step.\par}
  \bigskip\bigskip
  {\color{AtlasGoldLight}\pgfornament[width=26mm,color=AtlasGoldLight]{71}\par}
\end{frame}}

\end{document}
